\section{Discusión}

\subsection{Barreras Económicas y Culturales}
A pesar de los avances, persisten desafíos económicos significativos: muchos productores operan con márgenes ajustados, priorizando necesidades básicas sobre inversiones tecnológicas. La resistencia cultural al cambio, especialmente en poblaciones mayores de 50 años, requiere estrategias de sensibilización continua y demostración de beneficios tangibles.

\subsection{Desafíos de Conectividad Rural}
La brecha digital en el Huila limita el acceso a internet, afectando la sincronización de datos y la comunicación en tiempo real. Aunque el modo offline mitiga parcialmente este problema, se requieren inversiones en infraestructura comunitaria para una adopción masiva.

\subsection{Comparación con Literatura Relacionada}
Los resultados del portal se alinean con estudios sobre transformación digital en agricultura (Barrios-Ulloa et al., 2025; CEPAL, 2020), confirmando que la inclusión financiera y la capacitación son claves para la sostenibilidad. Sin embargo, el contexto colombiano presenta desafíos únicos de desigualdad que requieren enfoques híbridos offline/online.

\subsection{Implicaciones para Políticas Públicas}
El proyecto demuestra que la tecnología puede ser un motor de desarrollo rural, pero requiere alianzas entre gobierno, academia y sector privado. La promoción de Wi-Fi comunitario y subsidios para equipos podrían acelerar la adopción, contribuyendo a la soberanía alimentaria y reducción de pobreza.

\subsection{Limitaciones y Amenazas a la Validez}
Las pruebas piloto se limitaron a veredas específicas, por lo que la generalización requiere escalamiento gradual. Factores externos como variaciones climáticas pueden influir en los resultados económicos, sugiriendo la necesidad de evaluaciones longitudinales.

\subsection{Impacto del Cambio Climático}
El cambio climático nos trae cambios fuertes: temperaturas extremas, lluvias irregulares y nuevas plagas. Los agricultores necesitan herramientas para anticipar esos riesgos, planificar el riego, y saber qué semillas son más resistentes. Los datos climáticos evitan pérdidas, garantizan el abastecimiento y estabilizan la oferta en épocas difíciles. La capacidad de adaptación al clima es la base de un sistema agrícola sostenible.

Implementación en el proyecto (En modo amigo): El portal podría integrar el pronóstico climático por tu municipio, alertas de heladas o lluvias fuertes, y recomendaciones de siembra para cada estación. También podrás registrar el clima histórico de tu finca para planificar mejor tu año.

El campesino vive y respira con el clima. El cambio climático es una realidad que nos está dando muy duro, con heladas sorpresivas o aguaceros que llegan cuando no deben. Esto no es solo un problema técnico; es un golpe directo a la tranquilidad de la familia y a la posibilidad de asegurar el plato de comida.

El módulo climático del portal es un seguro de vida digital para tu cosecha. Te da la mano para que no trabajes a ciegas. Ya no tendrás que adivinar cuándo sembrar, porque tendrás el pronóstico de tu municipio. Te avisa si viene una helada fuerte, dándote tiempo para proteger tus plantas.

\subsection{Visión Futurista del Agro Colombiano}
El futuro del campo no es manual: estará lleno de sensores de humedad, drones de monitoreo, riego automático, inteligencia artificial (IA) y comercio 100\% digital. El productor dejará de depender del intermediario y se convertirá en un empresario agro-tecnológico.

La agricultura del futuro será inteligente, conectada y global. Colombia tiene todo para ser líder si invertimos en internet, educación rural y plataformas comerciales. La tecnología no reemplaza al campesino: lo dignifica y le da poder.

Implementación en el proyecto (En modo amigo): El Portal Agro Comercial del Huila es el primer paso para esa gran transición. Será la base para integrar en el futuro módulos de riego inteligente, monitoreo satelital y predicción de cosechas con IA, convirtiéndote en el productor del siglo XXI.

\subsection{Barreras Culturales y de Adopción}
La resistencia al cambio es un obstáculo común en comunidades rurales mayores de 50 años. El miedo a lo desconocido y la falta de familiaridad con tecnologías digitales frenan la adopción.

Estrategias como capacitaciones presenciales y demostraciones prácticas son esenciales para superar estas barreras. El portal debe ser intuitivo y amigable para ganar confianza.

\subsection{Impacto en la Economía Local}
La digitalización no solo beneficia al productor individual, sino que fortalece la economía regional. Al eliminar intermediarios, más ingresos quedan en las comunidades, impulsando desarrollo local.

Esto incluye creación de empleos en soporte técnico y logística, además de fomentar emprendimientos relacionados con la agricultura digital.

\subsection{Desafíos de Escalabilidad}
Escalar el portal a otras regiones requiere infraestructura robusta. La variabilidad en conectividad entre municipios plantea retos para una implementación uniforme.

Soluciones incluyen alianzas con proveedores de internet y modelos híbridos offline/online para asegurar accesibilidad universal.

\subsection{Comparación con Iniciativas Globales}
Comparado con proyectos en Europa o Asia, el portal colombiano se destaca por su enfoque en inclusión social. Mientras otros se centran en eficiencia, aquí se prioriza equidad y soberanía alimentaria.

Esto posiciona a Colombia como líder en agricultura digital inclusiva, inspirando modelos similares en países en desarrollo.